% Select, Compress, Reinvest -- ACL/ARR SUBMISSION VERSION.
% Derived from arxiv.tex on 2026-08-21. DO NOT edit this file directly for
% science changes -- make them in arxiv.tex and re-derive, or the two diverge.
% Difference from arxiv.tex is LENGTH ONLY: no claim, number, or caveat is
% dropped. Everything cut from the body is moved to the appendix, which ARR
% does not count. Sites are marked  % ACL-TRIM: <what/where it went>.
% ARR: switch to \usepackage[review]{acl} and drop the GitHub URL.
% Select, Compress, Reinvest: controlled visual-token allocation study.
%
% =====================================================================
% REVISION LOG -- 2026-08-21 post-review pass (5-seat panel)
% Snapshot before this pass: archive/arxiv.tex.pre-review-2026-08-21
% Each change site is marked inline with  % REV-2026-08-21 [Rn] ...
% Items marked TODO are NOT yet done and need a decision or a run.
%
% [R1] DONE. Tab:published rebuilt. Prior LDDR row cited the WRONG CELL:
%      +3.92 was Qwen2.5-VL-7B / 32F / full LDDR (stage1+2). The matched
%      cell is Qwen3-VL-8B / 8F / stage-1 LD (LDDR Table 1, p6 of
%      arXiv:2605.11477v1). Our reproduction is BELOW theirs on 8 of 9
%      cells, not above. All comparative verbs removed; reframed as a
%      cross-harness disagreement measurement. VERIFY against the
%      published version before camera-ready.
% [R2] DONE. Sec:scorer -- p-value for the OMP-vs-topk gap corrected from
%      .73 to .6752. The .73 was a two-arm test (OMP vs uniform across
%      scorers, uniform cancels); the sentence needs the four-arm
%      difference-in-differences. Point estimates coincide (+0.97) only
%      because top-k's marginal accuracy is identical under both scorers.
% [R3] DONE. Deleted a sentence duplicated verbatim (old lines 404/405).
% [R4] DONE. Cost claims qualified: every "cheaper"/"half the frames" now says
%      "in answerer visual tokens". Selection-stage compute (1 fps decode
%      + full-pool LongCLIP encode) is NOT counted anywhere in this paper.
% [R5] DONE. Scorer-invariance softened. "leaves it unchanged" -> "leaves the
%      ordering intact"; "and its significance" deleted from the
%      contribution; MDE/CI stated. Panel flagged MDE ~7.9pt > the 7.8pt
%      main effect, i.e. absence of evidence.
% [R6] DONE. Discordant counts (b,c) added for every Sec:scorer contrast.
% [R7] DONE. AKS marked as scorer-substituted; the star was previously applied
%      to FOCUS only, though AKS natively scores with CLIP+BLIP branches.
% [R8] DONE. Modality gap cited (liang2022modality); wording fixed from "orthogonal to
%      "the space of frame embeddings" to "the dominant directions of the
%      frame cone" (the original was false; a random vector carries ~54%
%      of its norm in a 412-frame span in 768-d).
% [R9] DROPPED 2026-08-29, deliberately. MDP3 and Q-Frame would each be a
%      from-scratch port with no published number at our budget to validate
%      against -- exactly the failure mode that produced the AKS padding bug
%      (see [AKS-FIX] below). The residual-saturation mechanism in Sec:why
%      covers unrun selectors better than two more unvalidated rows would.
%      Reopen only with a validation target in hand.
% [R10] DONE. VERIFIED against arXiv:2605.11477v1 p7 Table 2b: under one total
%      token budget, 1024 tok/frame 59.76, 512 tok/frame 61.03, 256 tok/frame
%      59.31 (LongVideoBench, Qwen2.5-VL-7B). Halving per-frame tokens doubles
%      frames, so this IS the reinvestment tradeoff, with an interior optimum.
%      Now cited and Claim B repositioned as an audited/paired/selector-agnostic
%      replication. What stays ours: measured token parity (0.996/0.984 vs their
%      asserted parity), paired McNemar, and the uniform 8->16 control showing
%      the effect does not require OMP.
% [R11] DONE 2026-08-29. Blind/text-only floor on LongVideoBench-600s:
%      .4369 (n=412) -- 43.7% answerable with NO frames. blind->uniform +10.9,
%      uniform->OMP +9.95. Per-benchmark floors for Video-MME and LVBench are
%      still NOT run; the reported floor is LVB-600s only.
% [R12] DONE 2026-08-29 (LVB-600s, n=412, all five selectors). The fused
%      query is a LEVEL SHIFT (+2.9..+3.9 for every rule except FOCUS*, -0.49),
%      not a re-ranking: the top three hold position. Table 1's ordering is not
%      an artifact of the scorer's query text. Run on LVB-600s only.
% [R13] DONE 2026-08-29. Full k=32 matrix run on LVB-600s (11 arms, n=412).
%      No selector-vs-selector difference is significant at either budget; every
%      selector beats uniform at k=8 but only AKS/LDDR-select/OMP do at k=32
%      (uniform gains +6.80 from k=8->32, OMP +1.70). The ordering is
%      unidentified at k=32, not reversed. LVB-600s only.
% =====================================================================

% SINGLE SOURCE OF TRUTH. Earlier versions (main.tex, main-old.tex) are in archive/.
% arXiv: \usepackage[preprint]{acl} (authors visible).
% ARR:   \usepackage[review]{acl}  (anonymous; also drop the GitHub URL in Sec. "Data and Code").
\documentclass[11pt]{article}
\usepackage[preprint]{acl}
\usepackage{times}
\usepackage{latexsym}
\usepackage{booktabs}
\usepackage{graphicx}
\usepackage{amsmath}
\usepackage{amssymb}
\usepackage{xcolor}
\usepackage{multirow}
\usepackage{url}

\graphicspath{{../figures/}}

\title{Select, Compress, Reinvest:\\A Controlled Study of Visual-Token Allocation in Long-Video MLLMs}

\author{Prakhar Khatri \\
  Independent Researcher \\
  \texttt{prakharkhatri123@gmail.com}}

\begin{document}
\maketitle

\begin{abstract}
% ACL-TRIM: abstract condensed 300 -> ~185 words. No claim removed; the
% audited/unaudited distinction and the scorer-swap power bound are kept.
An hour of video sampled once per second is 3{,}600 frames; a video language model can usually afford eight.
Choosing those eight well matters more than the field's default suggests: on LongVideoBench's hour-long bin, eight query-selected frames answer more questions correctly than sixteen uniformly spaced ones.
Establishing that comparison is hard, because published selectors change the frame scorer, the prompt, the resolution policy, and the answering model together.
We hold all four fixed and vary one decision at a time.
Under a single frozen LongCLIP scorer that sees only the question stem, we compare six training-free selection rules at eight frames on LongVideoBench, Video-MME, and LVBench.
Which rule is used matters more than whether it is query-aware at all---on Video-MME cosine top-$k$ gains 0.67 points over uniform while OMP gains 5.9---and the strong rule needs no new algorithm: classical Orthogonal Matching Pursuit~\citep{pati1993omp}, published in 1993, beats uniform sampling by 5.7, 5.9, and 11.8 points and stays within one point of a purpose-built modern selector.
A porting error we found and corrected in our own AKS baseline had made that row a duplicate of cosine top-$k$; correcting it changed 99.5\% of selected frames while moving LVBench accuracy by 0.07 points.
The uniform--top-$k$--OMP ordering also holds under SigLIP on one bin, within a test that bounds only large scorer effects, so the full six-rule ordering is not shown to be scorer-invariant.
Holding timestamps fixed, halving each frame's spatial budget moves accuracy by at most 0.44 points; spending the savings on more frames adds 2.36 points on LongVideoBench's long bins at verified equal cost.
Select relevant frames first, compress them second.
\end{abstract}

\section{Introduction}
% ACL-TRIM: Introduction condensed ~842 -> ~530 words. Every claim, number and
% caveat retained; the OMP-choice justification and the confound argument are
% compressed rather than cut. Full-length version in arxiv.tex.

A video language model cannot look at a whole video.
Decoded at one frame per second, a ten-minute clip is 600 images and an hour-long one is 3{,}600; a typical long-video system encodes eight.
Everything the model will know about the video passes through that handful, so the rule picking them is the first and tightest bottleneck, not a preprocessing detail.

The default rule is still uniform sampling, regardless of what was asked, and our clearest result is that this default is expensive: on hour-long LongVideoBench videos, eight frames chosen for relevance to the question beat sixteen uniformly spaced ones by 6.9 points ($p{=}.0011$), and the pattern holds one budget lower on ten-minute videos.
Half the frames, better answers: spending an \emph{answerer visual-token} budget on \emph{which} frames to keep buys more than spending it on \emph{how many}, before counting the selection stage's own cost.

Many systems already replace uniform sampling with query-aware retrieval, coverage, diversity, or learned evidence scores~\citep{tang2025aks,sun2025mdp3,zhu2025focus,wang2026evidential,peng2026qca}, and others jointly adapt frame choice and resolution~\citep{zhang2025qframe,chen2026lddr,wang2026dafs} or prune tokens inside the vision stack~\citep{vistallm2026,moprune2026}.
They report large end-to-end gains but leave a narrower question open: holding the scorer, frame budget, and answering model constant, which rule supplies the most useful evidence? Published tables cannot answer it, because those components move together.
A selector can look strong because it ships with a better encoder; a dynamic-resolution system can look strong because it fits in more frames rather than because it allocates resolution well.
Comparing two numbers from two papers compares two experiments, not two ideas.

This paper isolates the pieces.
We fix one frame scorer, one prompt boundary, one answering model, and one evaluation harness, then intervene on a single decision at a time.
\textbf{Select} changes which timestamps are chosen at fixed count and resolution; \textbf{compress} holds those timestamps and shrinks each frame's spatial budget; \textbf{reinvest} spends the recovered budget on more timestamps rather than sharper ones.
Each step is a paired comparison on the same questions, so its difference is attributable to that step alone.
A fourth comparison swaps the scorer itself, checking that the resulting ordering is not an artifact of the encoder we happened to freeze.

The rule at the centre of these interventions is deliberately old.
Orthogonal Matching Pursuit~\citep{pati1993omp} is a greedy sparse-approximation algorithm from 1993: it repeatedly takes the candidate most correlated with what the query still needs, then projects that direction away before choosing again.
We adopt it unmodified and untuned, which is a methodological choice rather than a concession.
An off-the-shelf rule has no hyperparameters for us to fit to these benchmarks, so a gap between it and a competing selector reflects the controlled comparison rather than unequal effort on our side of it.
It also sets an informative bar: how well the least sophisticated thing that could plausibly work does here measures how much recent purpose-built selectors gain from their subset rules as opposed to their scorers and pipelines.

The three interventions rank consistently across the settings we tested, but they were estimated in separate contrasts rather than one factorial design, so the ordering is a synthesis rather than an identified effect ranking.
Selection is the large lever, improving on uniform sampling by 5.7 to 11.8 points depending on the benchmark and beating uniform inputs that carry twice as many frames.
Compression is nearly free but useless alone: halving the per-frame budget at fixed timestamps changes accuracy by at most 0.44 points, with formal $\pm3$-point equivalence on pooled long videos.
Reinvestment converts that free budget into two to three points.
None of this is universal---the benefit depends on which model reads the frames and which benchmark asks the question, and OMP's later picks drift toward content that is visually novel but irrelevant---and we report those boundaries rather than smoothing them.

\paragraph{Contributions.}
\begin{itemize}
  \item \textbf{A matched comparison of six training-free selectors} under one scorer, one prompt boundary, one frame budget, and one answerer harness. OMP is a strong query-focused rule and stays within one point of LDDR's stage-1 selector.
  \item \textbf{Evidence that selection substitutes for frame count.} OMP significantly outperforms uniform sampling while using half as many frames, in two disjoint LongVideoBench duration bins (600 and 3600\,s) that share a benchmark, scorer, answerer, and implementation.
  \item \textbf{A decomposition of the fixed-budget allocation decision.} Roughly halving the per-frame spatial budget preserves accuracy, bounded by a two-one-sided-test interval rather than a bare null result, at a margin chosen post hoc; reinvesting the savings in additional keyframes then improves long-video QA.
  % REV-2026-08-21 [R5] deleted "and its significance"; added the power bound.
  \item \textbf{Evidence that the ranking survives a scorer swap.} Replacing LongCLIP with SigLIP changes 67--84\% of the selected frames yet leaves the selector ordering intact, within a test that bounds only scorer effects above roughly five points.
  \item \textbf{Evidence that aggregate accuracy conceals implementation error.} A padding branch in our own AKS port made that baseline select global top-$k$ at every $k<32$. Correcting it moved roughly 99.5\% of selected frames yet changed LVBench accuracy by 0.07 points. We report the cross-arm overlap check that caught it.
  \item \textbf{Explicit boundaries on these claims}, from paired tests, two answerer families, three benchmarks, a residual-geometry diagnostic, and a purposive failure audit. The boundaries are uneven: the scorer swap covers one bin, the equivalence result one pooled setting, and the audit is exploratory.
\end{itemize}

\section{Related Work}
% ACL-TRIM: Related Work condensed ~434 -> ~250 words. Every citation retained;
% the per-method one-line descriptions are compressed, and each paragraph keeps
% its gap statement, which is what the section is for.

\paragraph{Query-aware frame selection.}
Training-free selectors combine relevance, timeline coverage, and diversity in different proportions: AKS recursively partitions the timeline~\citep{tang2025aks}, FOCUS treats clips as bandit arms~\citep{zhu2025focus}, and MDP3, AdaRD-Key and Adaptive Greedy use DPP-based programs~\citep{sun2025mdp3,adardkey2025,adaptivegreedy2026}. Segment-based methods add per-segment quotas~\citep{peng2026qca} or MMR-refined event segments~\citep{chen2026efs}; others add uncertainty-triggered computation~\citep{kim2026request} or learned frame-level evidence~\citep{wang2026evidential}.
These differ not only in the subset rule but in scorer, candidate pool, training requirements, and frame budget, so reported gaps confound the rule with its pipeline. Our comparison removes that confound by giving every rule the same cached scores.

\paragraph{Selection with spatial adaptation.}
A second line couples frame choice to resolution~\citep{zhang2025qframe,chen2026lddr,wang2026dafs}. A joint gain has at least three possible sources---better timestamps, better resolution policy, or simply more temporal coverage---and none of these papers separates them, which is what our select/compress/reinvest sequence does.

\paragraph{Visual-token reduction.}
Token-level methods cut cost during or after encoding~\citep{adacodec2026,adapttoken2026,vistallm2026,moprune2026}. Concurrent work poses our question directly: AdaAlloc~\citep{an2026adalloc} asks how much of a fixed budget to spend on global context versus high-resolution local evidence; no preprint or code was public at the time of writing, so we describe it from its project-page abstract and draw no comparison against our numbers.
We ask a coarser but separable question upstream of all of them---before any token-level pruning, does an answerer gain more from sharper frames or from additional moments?

\section{Method: Three Controlled Interventions}
\label{sec:method}

% ACL-TRIM: protocol schematic moved to appendix (full-width float, ~0.5pp).
% It stays in the body in arxiv.tex, where it belongs.


\subsection{The experimental unit}

Every comparison here is paired: the unit is one benchmark question answered twice, once under each of two input policies, everything else held constant.
That counts the questions each policy wins rather than comparing two accuracy numbers which may differ for unrelated reasons, and measuring one decision at a time is what separates this from a comparison of complete systems whose components all differ at once.
The interventions differ only in what the two policies may change: \textbf{selection} changes timestamps at fixed frame count, scorer, and resolution; \textbf{compression} holds timestamps fixed and changes the per-frame spatial budget; \textbf{reinvestment} spends the saved budget on additional timestamps.
A fourth sits outside that sequence: replacing the scorer while holding rule, budget, and answerer fixed asks whether the resulting ordering is a property of the rules or of one encoder (\S\ref{sec:scorer}).

\subsection{One scorer, shared by every rule}

Videos are decoded at 1\,fps.
Each candidate frame and the question stem are encoded once with LongCLIP~\citep{zhang2024longclip}, and the resulting embeddings and stem similarities are cached.
Every selector then reads the same cache, so no rule benefits from a better encoder than another.

Answer options never enter selection.
That is a convention rather than a neutral choice, so we measure what it costs.
Scoring against the fused question-and-options string instead of the stem alone changes 41.9\% of the frames top-$k$ selects on LongVideoBench-600\,s and 53.3\% of those OMP selects; OMP is the more sensitive of the two because each pick is chosen to explain a residual component of the query vector, so a perturbation of that vector compounds through the greedy chain.
% REV-2026-08-22 [R15] Same refuted mechanism claim as arxiv.tex:228; see the
% revision-log entry there. Compressed here for the ACL page limit: the arXiv
% version reports the control arm's own contrasts (+2.18 p=.23, +1.70 p=.42).
On that bin, with the fused query, OMP reaches .6699 against .6311 for the stem (${+}3.88$ points, $p{=}.033$) and top-$k$ reaches .6408 against .6068 (${+}3.40$, $p{=}.087$).
A pre-registered replication on the 3600\,s bin agrees in sign and magnitude: .5691 against .5461 (${+}2.30$ points, $n{=}564$, $p{=}.18$), with the fused query moving roughly 69\% of the selected frames. Neither long bin resolves the shift alone, but two disjoint video sets agree, so we read the prompt boundary as a level shift on long video rather than a property of one bin.
Those are better frames, but not because of the options: a control that swaps in an unrelated video's options, matched on length, already reaches .6529---56\% of the gain---so the effect is not specific to the answer set, and neither contrast separates at this sample size.
We nonetheless score with the stem alone, following the official AKS implementation, which passes the question without candidates to its CLIP and BLIP branches~\citep{tang2025aks}; every selector here is therefore in a deliberately conservative configuration, and the gap above is the price of that choice.
The consequence for reading the literature is the substantive one: the text handed to the scorer is an uncontrolled axis worth close to four accuracy points, so a comparison inconsistent about it is measuring the prompt rather than the selector.

\subsection{Selection rules}

The matched comparison contains uniform sampling, cosine top-$k$, AKS~\citep{tang2025aks}, FOCUS$^\star$~\citep{zhu2025focus}, OMP, and LDDR-select~\citep{chen2026lddr}.
% REV-2026-08-21 [R7] the substitution is asymmetric and was previously marked on
% FOCUS only. AKS natively scores with CLIP and BLIP branches; LDDR-select alone
% runs on its native encoder, which is a standing advantage worth disclosing.
Every rule consumes our LongCLIP stem scores, but that is native only for LDDR-select: AKS and FOCUS both score with BLIP-ITM in their published form, so both are reproductions under a substituted scorer, and LDDR-select alone runs on its authors' encoder---an advantage to keep in mind when reading its position in Table~\ref{tab:t1}.
FOCUS$^\star$ carries a star for a second reason: it replays the published clip-bandit schedule against the dense LongCLIP vector rather than the original budgeted online ITM process (Appendix~\ref{app:focus}).
LDDR-select is the stage-1 Linear-DPP selector only, a published separately tabulated configuration rather than a truncation we invented.

Let $E\in\mathbb{R}^{N\times d}$ contain L2-normalized frame embeddings and let $q_0$ be the normalized stem embedding.
At each step OMP takes the frame most correlated with the current residual, then subtracts the span of everything already selected from the query:
\begin{equation}
\begin{aligned}
b_r &= \arg\max_{i\notin B_{r-1}} e_i^\top q_{r-1},\\
q_r &= q_0-\operatorname{Proj}_{\operatorname{span}\{e_j:j\in B_r\}}(q_0),
\end{aligned}
\label{eq:omp}
\end{equation}
where $B_r=B_{r-1}\cup\{b_r\}$.
In its original setting this is sparse reconstruction: the residual shrinks as selected atoms explain more of the target. Whether that survives in a contrastive image--text space is an empirical question, tested in \S\ref{sec:mechanism}; what the projection reliably does here is suppress directions already covered.

\subsection{Spatial budget and reinvestment}

For the fixed-timestamp intervention, frames are resized before reaching the model's processor, whose own internal resize is disabled for Qwen3-VL so that ours is the only one applied.
The main compressed arm targets a mean spatial fraction of about 0.53 (D@53); two controls test what the schedule's \emph{shape} contributes (Appendix~\ref{app:compression}).
The reinvestment arm instead selects $k{=}16$ timestamps under an approximately 50\% per-frame cap.
Because the comparison only means something if the two arms really cost the same, we audit tokens empirically rather than assuming the resize ratio transfers: the sixteen-frame arm consumes 0.996 and 0.984 times the tokens of the eight-frame full-resolution arm in the 600 and 3600\,s bins, so the winning arm is also the cheaper one.
Video-MME and LVBench use the same resize rule but were not separately audited, and we do not claim parity there.

\section{Evaluation Setup}
\label{sec:setup}

We evaluate on LongVideoBench~\citep{wu2024longvideobench} (the complete official validation set, $n{=}1337$), Video-MME~\citep{fu2025videomme} ($n{=}2700$), and LVBench~\citep{lvbench2025} ($n{=}1549$).
Tables abbreviate the first two as LVB and V-MME where column width requires it.
LongVideoBench is additionally stratified into 15, 60, 600, and 3600\,s duration bins, which is what makes the duration analysis in \S\ref{sec:selection} possible.
The primary answerer is Qwen3-VL-8B-Instruct~\citep{qwen3vl2025}, run through lmms-eval~\citep{zhang2024lmmseval} with greedy decoding, temperature zero, and subtitles disabled.
InternVL3-2B and InternVL3-8B~\citep{zhu2025internvl3} provide a second open-model family, and GPT-5-mini checks selected contrasts on the same saved timestamps.

% ACL-TRIM: table `tab:setup` moved to appendix.


\paragraph{Statistics.}
Accuracy comparisons use exact two-sided McNemar tests~\citep{mcnemar1947} on paired correctness outcomes, which is the appropriate test when the same question is answered under both policies.
Compression claims need more than a non-significant McNemar result, because failing to detect a difference is not evidence that none exists.
Where we claim two arms are interchangeable we therefore report two one-sided tests (TOST) with 90\% confidence intervals~\citep{schuirmann1987}, which bound the effect rather than merely failing to find it.
The $\pm2$, $\pm3$, and $\pm4$ percentage-point margins were chosen after seeing the runs rather than preregistered; we therefore report both the narrowest margin the data support and the adjacent margin they do not.
Tests are contrast-specific and uncorrected for multiplicity.

\section{Results}
\label{sec:results}

\subsection{Choosing frames beats having more of them}
\label{sec:selection}

The most direct evidence that selection matters is that it substitutes for frame count.
On 3600\,s LongVideoBench videos, OMP with eight frames reaches .5461 while uniform sampling with sixteen frames reaches .4770---a 6.9-point gain on half the input ($p{=}.0011$).
The same pattern holds one budget higher: on 600\,s videos, OMP with sixteen frames reaches .6578 against .6044 for uniform sampling with thirty-two, a 5.3-point gain at half the frame count ($p{=}.018$).
% REV-2026-08-21 [R4] "cheaper" scoped to answerer visual tokens. Selection-stage
% compute (1 fps decode + full-pool LongCLIP encode) is not counted anywhere here.
% ACL-TRIM: condensed; the selection-cost caveat is retained in full.
In both, the policy that is cheaper \emph{in answerer visual tokens} wins---a stronger statement than any equal-budget comparison can make, though the accounting covers the answerer only: producing the selected set also requires decoding at 1\,fps and encoding the whole pool with LongCLIP, a cost we do not measure and which a single-query-per-video deployment would not amortise.
Doubling the frame count is the obvious way to spend more compute on a long video, and in these two LongVideoBench bins it is the worse way.

\begin{table*}[t]
\centering
\small
\begin{tabular}{lccc}
\toprule
Method & LongVideoBench & Video-MME & LVBench \\
\midrule
Uniform & .5654 ($-$5.69) & .5637 ($-$5.85) & .3454 ($-$11.81) \\
Top-$k$ & .6028 ($-$1.95) & .5704 ($-$5.18) & .4319 ($-$3.16) \\
\textbf{OMP} & .6223 & \textbf{.6222} & .4635 \\
AKS & .5916 ($-$3.07) & .6059 ($-$1.63) & .4287 ($-$3.48) \\
FOCUS$^\star$ & .5819 ($-$4.04) & .5578 ($-$6.44) & .3983 ($-$6.52) \\
LDDR-select & \textbf{.6320} (${+}0.97$) & .6193 ($-$0.29) & \textbf{.4693} (${+}0.58$) \\
\bottomrule
\end{tabular}
\caption{Matched selector comparison at $k{=}8$ with Qwen3-VL-8B. Cells are accuracy; parentheses give the difference from OMP in percentage points. FOCUS$^\star$ is the controlled schedule replay described in Appendix~\ref{app:focus}; LDDR-select is stage~1 only.}
\label{tab:t1}
\end{table*}

Table~\ref{tab:t1} places that result in the wider comparison, with all six rules reading the same cached scores.
OMP improves on uniform sampling by 5.69 points on LongVideoBench, 5.85 on Video-MME, and 11.81 on LVBench, and the paired evidence is decisive on Video-MME ($p{=}5.1\times10^{-10}$), LVBench ($p{=}9.5\times10^{-17}$), and the LongVideoBench 3600\,s bin ($p{=}5.9\times10^{-4}$).
LDDR-select is the one rule that keeps pace, staying within a point of OMP on all three benchmarks with no detectable paired difference on Video-MME or LVBench ($p{=}.70$ and $.57$).

% ACL-TRIM: figure `fig:longvideo` moved to appendix.


The duration split explains when any of this matters.
On 15\,s clips, uniform sampling, top-$k$, and OMP return exactly the same accuracy (.7249): eight frames out of fifteen candidates cover the video regardless of how they are picked, so there is nothing for a selector to do.
The gains appear once the candidate pool outgrows the budget---7.8 points at 600\,s ($p{=}.0022$) and 7.5 points at 3600\,s ($p{=}5.9\times10^{-4}$).
% REV-2026-08-30 [SCOPE] the caption previously said the gain "grows with duration",
% which the paper's own numbers contradict: 7.8 at 600s vs 7.5 at 3600s.
Note that the gain does not keep growing after that: 7.5 at 3600\,s is if anything slightly below the 7.8 at 600\,s, so this is a threshold that switches on when the pool exceeds the budget, not a gradient in duration.
The duration bins also differ in content and question composition, so this is an observational contrast between bins; isolating pool size from duration would require a controlled truncation experiment that we did not run.

Which \emph{kind} of rule is needed is more benchmark-dependent than the headline suggests.
% REV-2026-08-30 [AKS-FIX] AKS moved to the other side of this sentence: the
% corrected rule separates from uniform on Video-MME at p=4.0e-6 (n=2700).
On Video-MME, top-$k$ and FOCUS$^\star$ do not separate detectably from uniform sampling, while OMP, LDDR-select, and AKS do---AKS by 4.22 points ($p{=}4.0\times10^{-6}$); on LongVideoBench, OMP's advantage over plain top-$k$ is small and often undetectable within individual bins.
The supported reading is narrow and worth stating precisely: combining query relevance with suppression of already-covered directions is consistently strong in this harness, whereas the value of any individual simpler heuristic does not transfer across benchmarks on its own.
Whether that ordering is a property of the rules or of LongCLIP is the subject of \S\ref{sec:scorer}.

% REV-2026-08-21 [R5] heading asserted a null; now states what was tested.
\subsection{The ranking survives a scorer swap}
\label{sec:scorer}

Every number in Table~\ref{tab:t1} comes from one frozen scorer, leaving open whether the ordering reflects the selection rules or LongCLIP's geometry.
We replace LongCLIP with SigLIP-so400m~\citep{zhai2023siglip}, holding prompt boundary, budget, answerer and harness fixed, re-running both arms together so every contrast is paired on identical items.
The swap is substantial rather than cosmetic: on the 600\,s bin the scorers agree on only 1.32 of eight frames for OMP and 2.62 for top-$k$, so 67 to 84\% of what the answerer sees changes.

The ordering survives (Table~\ref{tab:scorer}): uniform is beaten by top-$k$ and top-$k$ by OMP under both scorers, with OMP over uniform significant in each case ($p{=}.0024$ under LongCLIP, $p{=}.00026$ under SigLIP).
Neither selector's gain depends detectably on which scorer produced it, but with roughly 134 discordant pairs the 90\% interval on the OMP scorer effect spans about $[-3.6,+5.6]$ points, so this bounds large scorer effects rather than establishing their absence.
Appendix~\ref{app:scorer} gives the 60\,s bin, the discordant counts, and an interaction test that rejects the tempting reading that OMP suits SigLIP better.

The claim is bounded: one alternative scorer on two duration bins does not establish scorer-invariance, and both encoders share a contrastive image--text objective, so an ITM head or attention-based scorer could still reorder the table.
What it weakens, without ruling out, is the worry that Table~\ref{tab:t1} is an artifact of LongCLIP: the ordering holds for the three rules tested when most selected frames are replaced.

\subsection{Half the pixels change almost nothing}
\label{sec:compression}

Table~\ref{tab:t2} holds the selected timestamps fixed and changes only how many pixels each frame is worth: across the three Qwen benchmark aggregates the compressed arm shifts accuracy by at most 0.44 points in either direction.
A uniform-timestamp control moves 0.31 points ($p{=}.784$), consistent with the stability belonging to the compression rather than to OMP-selected frames, though a non-significant control is not itself an equivalence result.

% ACL-TRIM: tab:t2 moved to Appendix~\ref{app:published}; every cell it carries
% is already stated numerically in the body prose above and below.


Small differences are easy to over-read, so we bound one of them properly.
On the pooled LongVideoBench long bins ($n{=}976$), the compressed OMP arm is 0.72 points \emph{higher} than full resolution, with a 90\% confidence interval of $[-0.60,+2.03]$: inside a $\pm3$-point margin (TOST $p{=}.0022$) but not $\pm2$ ($p{=}.054$). The margin was chosen after seeing the interval, so this bounds the effect descriptively rather than confirming a pre-specified one.
Appendix~\ref{app:compression} gives the uniform-timestamp control, a stronger InternVL3-8B result under a more aggressive intervention, and the controls separating the compression \emph{level} from its allocation \emph{shape}; the short version is that coarse compression is close to free, a naive priority-proportional rule adds nothing, and a better-designed importance signal can add something.

\subsection{Spending the savings on more frames}
\label{sec:reinvest}

% REV-2026-08-21 [R10] Claim B has a published precedent. LDDR Table 2b (verified,
% arXiv:2605.11477v1 p7) reports fixed-resolution ablations under one total token
% budget: 1024 tok/frame 59.76, 512 tok/frame 61.03, 256 tok/frame 59.31 on
% LongVideoBench with Qwen2.5-VL-7B. Halving per-frame tokens doubles the frames,
% so that IS the reinvestment tradeoff, with an interior optimum. Repositioned
% below as an audited, paired, selector-agnostic replication rather than a new
% observation. Do not restate B as novel.
% ACL-TRIM: condensed; the precedent concession is kept intact -- do not cut it.
The direction is not new: \citet{chen2026lddr} report fixed-resolution ablations under one token budget where 1024 tokens per frame reach 59.76 on LongVideoBench, 512 reach 61.03 and 256 fall back to 59.31, so halving the per-frame budget doubles the frames that fit and gains 1.27 points before declining.
That ablation does not establish whether the parity it assumes holds, whether the difference exceeds noise, or whether the effect needs their selector; we address those three, making this an audited paired replication rather than a new observation.
Table~\ref{tab:t2}b spends it on temporal coverage, comparing eight full-resolution frames against sixteen compressed ones at equal or lower measured token cost: LongVideoBench improves 2.24 points overall and 2.36 on the pooled 600 and 3600\,s bins ($p{=}.0346$), LVBench 3.04 ($p{=}.0009$), Video-MME 1.56, and GPT-5-mini on LVBench 2.39 ($p{=}.039$).

% ACL-TRIM: merged into Table~\ref{tab:t2}.


% ACL-TRIM: figure `fig:claim-b` moved to appendix.


% ACL-TRIM: condensed; all four cells and every statistic retained.
We cannot show that this requires OMP.
All four cells of the selector-by-budget design---uniform and OMP, each at eight full-resolution and sixteen compressed frames---were run in one environment.
Uniform gains ${+}1.13$ points ($p{=}.3673$) on the eight-to-sixteen move and OMP ${+}2.56$; the interaction is 1.43 points and not significant ($p{=}.2614$), splitting by duration (${+}0.00$ at 600\,s, ${+}2.48$ at 3600\,s, neither individually significant).
With 203 discordant pairs this excludes only large effects, so the defensible claim is that temporal reinvestment works with OMP, points the same way under uniform sampling, and remains open on selector-specificity at roughly the one-point scale.

\subsection{What the later picks actually add}
\label{sec:mechanism}

A residual trace explains why late picks still help without reconstructing the query: the residual norm falls only from 0.972 to 0.967 across eight picks, while the residual's correlation with the selected frame collapses from .233 to .004, so the query-directed signal is exhausted early and later picks add coverage rather than new query direction.
This is a hypothesis consistent with the geometry, not a demonstrated mechanism; nothing here intervenes on residual geometry.
Appendix~\ref{app:mechanism} gives the full trace, the modality-gap account, and an exploratory audit of 93 OMP failures whose most frequent tag is off-topic novelty.

\subsection{The same frames do not help every model equally}
\label{sec:transfer}

Table~\ref{tab:internvl} hands identical OMP and uniform timestamps to InternVL3, changing the answerer and nothing else.
The selection gain replicates on LongVideoBench and LVBench at both sizes. Video-MME does not cooperate: the 2B model benefits at every duration while the 8B model benefits only on short videos, flat on medium and long despite demonstrably different frames.
Capacity alone does not account for it, since the same 8B model gains 10.14 points on LVBench; a capacity-by-benchmark interaction remains possible and one model family cannot separate the two.

% ACL-TRIM: table `tab:internvl` moved to appendix.


The allocation results transfer more smoothly.
On pooled LongVideoBench long videos, InternVL3-8B holds its accuracy through a threefold tile reduction ($-0.10$ points, $p{=}1.000$) and gains 3.38 points when the same total tile budget is spread from 8 uniformly sampled frames to 24 ($p{=}.0099$), an effect concentrated in the 600\,s bin.
Both halves of the story therefore survive a change of answerer but their magnitudes do not: these are properties of an answerer--benchmark pair, not constants of a timestamp set.

\section{Discussion}
\label{sec:discussion}

% ACL-TRIM: first two Discussion paragraphs condensed ~220 -> ~140 words.
\paragraph{The three decisions are not equally worth optimizing.}
Choosing frames is the large lever, sometimes beating uniform inputs carrying twice as many frames.
Compressing those frames is close to free but buys little accuracy alone, so its accuracy value lies in what the savings buy---temporal coverage returns a further two to three points---though compression may still pay for latency or memory when the tokens are not reinvested.
Fix selection first, treat resolution as the slack variable, and reinvest rather than pocket the difference.

\paragraph{Aggregate accuracy concealed a baseline that was reproducing top-$k$.}
Our own first AKS port ended in a padding branch absent from the released code: at $k{=}8$ its recursion allotted every leaf zero frames and the pad filled the selection from global cosine top-$k$, making that row a duplicate of the top-$k$ row on all three benchmarks.
Correcting it changes roughly 99.5\% of the selected frames, yet aggregate accuracy moves by $-1.05$ points on LongVideoBench and $+0.07$ on LVBench ($+2.74$ on Video-MME).
A single-number comparison cannot detect that, and neither can a per-arm check asking only whether each row looks plausible; what detected it was comparing arms against \emph{each other}, since rules optimising different objectives should not select the same frames (Appendix~\ref{app:published}).

\paragraph{A 1993 algorithm is competitive with purpose-built selectors.}
Orthogonal Matching Pursuit predates every method it is compared against by three decades, has no video-specific machinery and no tuned hyperparameters, and still leads AKS, FOCUS$^\star$, and top-$k$ on all three benchmarks while staying within a point of LDDR's stage-1 selector.
That is not evidence that recent selectors are poorly designed; it is evidence about where their gains live.
% REV-2026-08-21 [R14] previous version claimed the between-rule room was
% "considerably smaller than the distance from uniform to any query-aware rule".
% That FAILS on Video-MME: top-k gains 0.67 over uniform while OMP gains 5.85, so
% the between-rule spread (5.18 excluding FOCUS*) exceeds it. The fact was already
% reported in Sec:selection and simply not propagated here. Now bounded.
% REV-2026-08-30 [AKS-FIX] Bound restated again. With the corrected AKS row the
% between-rule span is 4.04 on LongVideoBench against a 3.74 uniform->top-k
% distance, so the claim now survives on LVBench ONLY (4.06 vs 8.65). Video-MME
% was already broken (R14). Do not repair this by dropping AKS from the span.
Once the scorer, prompt boundary, frame budget, and answerer are held fixed, the room between subset-selection rules is smaller than the uniform-to-top-$k$ distance on LVBench alone (4.06 against 8.65); on LongVideoBench the two are comparable (4.04 against 3.74) and on Video-MME the between-rule spread dwarfs it (5.18 against 0.67).
Video-MME is sharpest: top-$k$ gains 0.67 points over uniform while OMP, LDDR-select, and AKS each gain more than four, so \emph{which} rule is used matters far more than being query-aware at all.
The defensible generalization is therefore not ``subset rules barely matter''; and since we do not apportion variance between scorer, pipeline, implementation, and parsing, this compares two spans rather than attributing effects to components.
The corollary is concrete: a selector comparison is only informative if the scoring stage is shared, and a paper reporting a new rule should include a classical pursuit or greedy baseline on its own scores before claiming the rule is what helped.

% REV-2026-08-21 [R1] Table rebuilt from the MATCHED published cell.
% Previous row used LDDR's Qwen2.5-VL-7B / 32F / full-LDDR number (+3.92) and
% claimed we "exceeded" it. That cell differs from ours on model, budget, AND
% method variant. The matched cell (Qwen3-VL-8B, #F=8, stage-1 LD, LongCLIP;
% arXiv:2605.11477v1 Table 1, p6) shows our reproduction BELOW theirs on 8 of 9
% cells. All comparative verbs removed. Re-verify against the published version.
% ACL-TRIM: table `tab:published` moved to appendix.


% ACL-TRIM: the two published-vs-matched paragraphs moved to appendix; one
% pointer sentence kept here so the Discussion still records the finding.
\paragraph{Two matched harnesses disagree by as much as the effects they measure.}
Appendix~\ref{app:published} sets our gains beside the one prior study that standardises its encoder, at a cell matched on answerer family, budget, and encoder.
The two harnesses differ by 0.4 to 3.8 points on identical rules---a range covering most of the gaps published selector comparisons are built on---which is why every claim here is a paired contrast inside one harness rather than a comparison against a published figure.

% ACL-TRIM: final three Discussion paragraphs condensed ~330 -> ~200 words.
\paragraph{Joint allocation systems deserve decomposed reporting.}
This sharpens how to read systems such as Q-Frame, LDDR, and DAFS~\citep{zhang2025qframe,chen2026lddr,wang2026dafs}, whose end-to-end gains combine several causal routes.
Our controls find no value in a priority-proportional resolution rule, a positive effect from reconstructed LDDR-style importance, and the clearest downstream gain from temporal reinvestment---three answers a single end-to-end number blurs into one.
Methods would be easier to build on if they reported separate timestamp, resolution, and frame-count interventions alongside the joint result.

\paragraph{Conclusions should span durations and answerers.}
Selection is a property of the answerer--benchmark pair: InternVL3-8B is nearly insensitive to OMP on medium and long Video-MME yet gains over ten points on LVBench, and OMP's advantage is exactly zero at 15\,s.
A claim measured on one benchmark at one duration with one answerer is not safe to generalize; our own results would have supported at least three different headlines depending on which cell we reported.
The geometry points the same way: with the text query nearly orthogonal to the dominant directions of the frame-embedding cone~\citep{liang2022modality}, what remains useful is relevance plus decorrelation, and relevance floors or scorers with a smaller modality gap are the natural next interventions.

\section{Conclusion}
\label{sec:conclusion}

In our matched harness, which frames reach the answerer is the largest of the three allocation decisions we varied: eight well-chosen frames beat sixteen uniformly spaced ones in LongVideoBench's hour-long bin, and a greedy pursuit algorithm from 1993 delivers that, outperforming uniform sampling by 5.7 to 11.8 points and staying within a point of a purpose-built modern selector.
Roughly half the per-frame spatial budget can then be removed with little observed change, and reinvesting those tokens in additional keyframes returns a further two to three points.
How far this generalises beyond one harness is what the evaluation-sensitivity results bound: a porting error in our own baseline moved 99.5\% of selected frames while changing one benchmark's accuracy by 0.07 points, and two matched harnesses running the same published rules disagree by as much as the effects they measure.

\section{Limitations}
\label{sec:limitations}

\paragraph{Two selectors are controlled reproductions, not full systems.}
FOCUS$^\star$ replays the published selection schedule against dense LongCLIP scores instead of reproducing the original budgeted ITM scorer, so its row tests the schedule under a matched scorer and not the FOCUS pipeline's accuracy or efficiency.
LDDR-select is stage 1 only, and the stage-2 allocation contrast is reconstructed from the paper rather than run from the authors' code.
Neither row should be read as a reproduction result for the original method.

\paragraph{Equivalence margins were chosen after seeing the data.}
The Qwen result covers pooled LongVideoBench 600 and 3600\,s examples within a $\pm3$-point margin and misses $\pm2$. Because the margin was selected after inspecting the interval rather than fixed in advance, it answers ``which margin does this sample clear'' rather than ``does the effect fall inside an independently justified bound''. We therefore report the interval and both margins and treat the result as descriptive; a confirmatory equivalence claim would need a pre-specified margin and a fresh replication. The tests are also uncorrected for multiplicity.
Video-MME compression has no paired test, and the small LVBench differences are descriptive rather than equivalence evidence.
The selector-by-budget interaction is underpowered for effects near one point, which is the scale at which it would matter.

\paragraph{Mechanism evidence is encoder-specific.}
The residual geometry and Gram statistics come from frozen LongCLIP stem embeddings, mostly on LongVideoBench-600\,s.
The selector ordering survives a swap to SigLIP (\S\ref{sec:scorer}), but both encoders share a contrastive image--text objective; BLIP-ITM heads, MLLM-attention scorers, and subtitle-aware selection remain untested, and the modality-gap explanation in \S\ref{sec:mechanism} is a property of this embedding space rather than of pursuit selection in general.
The model-assisted failure audit examined selected rather than randomly sampled failures, so its counts are exploratory and cannot estimate population rates.

\paragraph{Some secondary results were measured in a second compute environment.}
Tables~\ref{tab:t1} and~\ref{tab:t2} and every claim drawn from them come from a single stack: RTX PRO 4500 GPUs with PyTorch 2.11/CUDA 13.0.
The selector-by-budget interaction in \S\ref{sec:reinvest}, the allocation-shape contrast in \S\ref{sec:compression}, and the variant sweep in Appendix~\ref{app:negatives} were measured on L40S GPUs with PyTorch 2.6.
In each case every arm of the contrast---including a re-run of the full-resolution baseline---was executed inside that second environment, so no reported difference is computed across the two stacks.
Replicating the baseline in both also lets us quantify the drift rather than assume it is small: the uniform eight-frame full-resolution arm reads 0.36 points higher on the second stack in the 3600\,s bin and is identical in the 600\,s bin, while the corresponding OMP arm reads 0.71 points lower.
Absolute accuracies from these sections are therefore not comparable to Tables~\ref{tab:t1} and~\ref{tab:t2} and we never place them in the same table; the paired differences within each section are unaffected.

\paragraph{Evaluation scope is incomplete.}
Subtitles are disabled throughout, which leaves subtitle-anchored questions partly inaccessible to every visual selector and depresses all arms together.
GPT-5-mini compression uses a pixel ratio rather than Qwen token accounting, and some GPT selector results survive only as aggregates rather than per-item records.

\section{Ethical Considerations}
\label{sec:ethics}

The experiments use existing public video-QA benchmarks and frozen open or API models.
No new human subjects, personal data, videos, or annotations were collected.
Question stems, but not answer options, enter the selector.
Benchmark videos remain governed by their original licenses and are not redistributed with the released artifacts.

\section{Data and Code Availability}

Evaluation code, selector implementations, analysis scripts, per-item predictions permitted by the benchmark licenses, and aggregate result tables are available at \url{https://github.com/codeprakhar25/omp-keyframe-sampling}.
The release includes the scorer-swap pipeline of \S\ref{sec:scorer}, the fused-query arms of \S\ref{sec:method}, and the script that builds Table~\ref{tab:published} from the cited papers' reported numbers.
The release excludes benchmark video files and API credentials.

\section{Author Contributions}

Prakhar Khatri: Conceptualization, Methodology, Software, Validation, Formal Analysis, Investigation, Data Curation, Visualization, Writing--Original Draft, and Writing--Review \& Editing.

\section{Funding}

This research received no external funding.

\section{Competing Interests}

The author declares no competing interests.

\section{AI Assistance Disclosure}

Generative AI tools assisted code development, manuscript editing, and an exploratory visual failure audit.
The author directed the analyses, verified the reported numerical results against saved artifacts, adjudicated the scientific claims, and takes responsibility for the final manuscript.

\bibliography{references}

\appendix



\section{Relocated tables and figures}
\label{app:published}
% ACL-TRIM: these are body content in the arXiv version (arxiv.tex); they are
% here only to meet the ARR page limit. Nothing is altered.
\begin{table}[h]
\centering
\small
\resizebox{\columnwidth}{!}{%
\begin{tabular}{lcccccc}
\toprule
& \multicolumn{2}{c}{Video-MME} & \multicolumn{2}{c}{LongVideoBench} & \multicolumn{2}{c}{LVBench} \\
\cmidrule(lr){2-3}\cmidrule(lr){4-5}\cmidrule(lr){6-7}
Rule & pub. & ours & pub. & ours & pub. & ours \\
\midrule
\emph{uniform (absolute)} & 57.19 & 56.37 & 54.53 & 56.54 & 28.08 & 34.54 \\
\midrule
AKS$^\dagger$ & ${+}4.29$ & ${+}4.22$ & ${+}5.01$ & ${+}2.62$ & ${+}12.07$ & ${+}8.33$ \\
FOCUS$^{\star\dagger}$ & ${+}1.70$ & $-0.59$ & ${+}5.00$ & ${+}1.65$ & ${+}3.49$ & ${+}5.29$ \\
LDDR-select & ${+}5.96$ & ${+}5.56$ & ${+}8.90$ & ${+}6.66$ & ${+}14.98$ & ${+}12.39$ \\
OMP & --- & ${+}5.85$ & --- & ${+}5.69$ & --- & ${+}11.81$ \\
\bottomrule
\end{tabular}}
\caption{Two matched-condition harnesses, same rules. Published columns are \citet{chen2026lddr}, Table~1, Qwen3-VL-8B at eight frames with a LongCLIP encoder---the same answerer family, budget, and encoder as ours---where LDDR-select is their stage-1 row. The first line gives absolute uniform-sampling accuracy; all other cells are gains over the uniform row of the same column. $^\dagger$ marks rules we run on a substituted scorer (both score with BLIP-ITM natively); $^\star$ marks the FOCUS schedule replay of Appendix~\ref{app:focus}. OMP has no published counterpart. MDP3 and Q-Frame appear in their table and we did not run them.}
\label{tab:published}
\end{table}

\begin{table}[h]
\centering
\small
\begin{tabular}{lccc}
\toprule
Selector & LongCLIP & SigLIP & SigLIP $-$ LongCLIP \\
\midrule
Uniform & \multicolumn{2}{c}{.5534} & --- \\
Top-$k$ & .6068 & .6068 & ${+}0.00$ ($p{=}1.00$) \\
OMP & .6311 & .6408 & ${+}0.97$ ($p{=}.73$) \\
\bottomrule
\end{tabular}
\caption{Scorer swap on LongVideoBench-600\,s ($n{=}412$), all arms run together and paired on identical items. Uniform sampling is scorer-independent and appears once. The final column is a paired test of whether each selector's gain depends on the scorer.}
\label{tab:scorer}
\end{table}

\begin{table}[h]
\centering
\small
\resizebox{\columnwidth}{!}{%
\begin{tabular}{lrrr}
\toprule
Benchmark & $n$ & IVL3-2B $\Delta$ & IVL3-8B $\Delta$ \\
\midrule
Video-MME short & 900 & ${+}6.00$ ($p{=}1.2\times10^{-4}$) & ${+}5.44$ ($p{=}1.9\times10^{-4}$) \\
Video-MME medium & 900 & ${+}3.33$ ($p{=}.046$) & ${+}0.33$ ($p{=}.891$) \\
Video-MME long & 900 & ${+}4.22$ ($p{=}.0062$) & ${+}0.11$ ($p{=}1.000$) \\
LVBench & 1549 & ${+}9.04$ ($p{=}7.2\times10^{-11}$) & ${+}10.14$ ($p{=}1.2\times10^{-13}$) \\
LongVideoBench & 1337 & ${+}4.71$ ($p{=}1.6\times10^{-4}$) & ${+}6.51$ ($p{<}10^{-4}$) \\
\bottomrule
\end{tabular}}
\caption{OMP minus uniform accuracy on identical timestamps and tile budgets with InternVL3. Deltas are percentage points.}
\label{tab:internvl}
\end{table}

\begin{figure*}[h]
\centering
\includegraphics[width=\textwidth]{fig1_residual_inert.pdf}
\caption{LongCLIP geometry on LongVideoBench-600\,s ($n{=}411$).
(a) Residual correlation vanishes while residual norm changes little; the line connects logged picks 1, 8, and 16 and is not a dense trace.
(b) Within-video frames occupy a coherent cone, while the text query is nearly orthogonal to the frame dictionary. In this space, OMP mainly suppresses directions already represented by earlier selections.}
\label{fig:residual}
\end{figure*}

\begin{figure}[h]
\centering
\includegraphics[width=\columnwidth]{fig2_long_video_effect.pdf}
\caption{OMP's gain over uniform sampling appears only once the 1\,fps candidate pool outgrows the eight-frame budget. At 15\,s, uniform, top-$k$, and OMP have identical accuracy (.7249); the gain is 7.8 points at 600\,s and 7.5 at 3600\,s. The effect switches on with pool size rather than growing with duration.}
\label{fig:longvideo}
\end{figure}

\begin{figure}[h]
\centering
\includegraphics[width=\columnwidth]{fig_claim_a.pdf}
\caption{LongVideoBench 3600\,s accuracy on fixed OMP-8 timestamps across spatial budgets. The flat 50\% control and the residual-proportional schedule are indistinguishable, so OMP rank alone does not provide a useful resolution policy.}
\label{fig:claim-a}
\end{figure}

\paragraph{Two matched harnesses disagree by as much as the effects they measure.}
LDDR is the one prior study that standardises every baseline to a single encoder, and it reports a cell matched to ours on answerer family, frame budget, and encoder.
Table~\ref{tab:published} places that cell beside our own.
Our gains are smaller in eight of the nine comparable cells, by 0.07 to 3.74 points.
We do not read this as either harness being wrong.
The uniform baselines themselves differ---most starkly on LVBench, 28.08 against our 34.54---so the two harnesses are not measuring from the same floor, and a lower floor leaves more room for a selector to gain.
Prompt template, candidate-pool construction, decoding path, and answer parsing all differ and none of them is specified tightly enough in either paper for a reader to reconcile the two.
% REV-2026-08-30 [AKS-FIX] Both halves of this sentence inverted once AKS was
% corrected: AKS is now closest on Video-MME (0.07) and LDDR on LongVideoBench.
We also find no clean pattern in the residuals, and in particular no advantage to keeping a rule's native encoder: on Video-MME the closest reproduction is AKS$^\dagger$, within 0.07 points of the published gain despite a substituted scorer, while LDDR-select---the one rule we run on its authors' own encoder---reproduces closest on LongVideoBench (2.24) and LVBench (2.59).

\paragraph{A baseline port that silently reproduced top-$k$.}
Our first implementation of AKS ended in a padding branch that is absent from the authors' released code.
Its recursive partition gives each leaf at depth $d$ a quota of $\lfloor k/2^{d}\rfloor$ frames, so at $k{=}8$ with the published depth of five every leaf is allotted zero, the recursion returns nothing, and our pad filled the entire selection from global cosine top-$k$.
For every benchmark in an earlier draft of this paper, the AKS row was a duplicate of the top-$k$ row.
We caught it because those two rows were the closest pair in Table~\ref{tab:t1} on all three benchmarks; a pairwise overlap check over selected frame \emph{indices} then showed 98.5\% agreement.
What is worth reporting is not the bug but how completely aggregate accuracy hid it.
Correcting the depth changes roughly 99.5\% of the selected frames on every benchmark, yet the aggregate moves by only $-1.05$ points on LongVideoBench and $+0.07$ on LVBench, while moving $+2.74$ on Video-MME.
On LVBench a reader comparing accuracies alone would see two-hundredths of a point between a faithful implementation and one selecting entirely different frames for 99.5\% of questions.
A single-number comparison cannot detect this class of error, and neither can a per-arm sanity check that only asks whether each row looks individually plausible.
What did detect it was comparing arms against \emph{each other}: if two rules that optimise different objectives select nearly the same frames, at least one is not implementing its objective.
We now run that overlap check by construction, and we recommend it wherever a table's rows are meant to be distinct methods.
The corrected row also reproduces the published Video-MME gain to within 0.07 points against a 2.81-point gap under the bug, which is evidence both that the correction is right and that this harness can match a published number when the port is faithful.

The disagreement is itself the measurement.
Two carefully controlled harnesses, running the same published rules at the same budget with the same encoder, differ by as little as 0.07 and as much as 3.74 points---a range that covers most of the gaps published selector comparisons are built on.
That bounds how much weight any cross-paper selector number can carry, including the numbers this paper's own argument targets, and it is why every claim we make is a paired contrast inside one harness rather than a comparison against a published figure.
It also means Table~\ref{tab:published} should be read as a measurement of harness disagreement, not as a reproduction verdict on any of these methods.

\begin{figure}[h]
\centering
\includegraphics[width=0.80\columnwidth]{fig_claim_b.pdf}
\caption{LongVideoBench 600\,s and 3600\,s pooled: sixteen compressed OMP frames outperform eight full-resolution OMP frames at a measured token ratio no greater than one.}
\label{fig:claim-b}
\end{figure}

\begin{table}[h]
\centering
\small
\begin{tabular}{@{}lp{0.58\columnwidth}@{}}
\toprule
Benchmarks & LongVideoBench (1337), Video-MME (2700), LVBench (1549) \\
Primary answerer & Qwen3-VL-8B, greedy decoding, subtitles off \\
Transfer checks & InternVL3 2B/8B; GPT-5-mini \\
Shared scorer & LongCLIP, question stem only, 1\,fps pool \\
Scorer control & SigLIP-so400m, identical protocol (\S\ref{sec:scorer}) \\
Frame budgets & $k{=}8$ for selection; $k{=}16$ for reinvestment \\
\bottomrule
\end{tabular}
\caption{Evaluation grid. Main claims concern paired accuracy differences under matched inputs, not absolute leaderboard position.}
\label{tab:setup}
\end{table}

\begin{figure*}[h]
\centering
\includegraphics[width=\textwidth]{fig_protocol.pdf}
\caption{Controlled protocol.
(a) A 1\,fps candidate pool is encoded once with LongCLIP using only the question stem; every $k{=}8$ selector consumes the same cached embeddings and scores.
(b) Selected timestamps are then frozen while spatial budget is reduced.
(c) The saved visual tokens are reinvested by comparing eight full-resolution frames with sixteen compressed frames at equal or lower measured token cost.
FOCUS$^\star$ denotes a replay of the published selection schedule on dense LongCLIP scores, not its original budgeted ITM scorer.}
\label{fig:protocol}
\end{figure*}

% ACL-TRIM: tab:t2 and tab:t3 merged into one two-panel float. No cell changed;
% this removes one caption/header block. They are separate tables in arxiv.tex.
\begin{table}[t]
\centering
\small
\resizebox{\columnwidth}{!}{%
\begin{tabular}{lccc}
\toprule
\multicolumn{4}{l}{\emph{(a) Fixed timestamps, halved spatial budget}} \\
Answerer / benchmark & Full & Compressed & $\Delta$ (pp) \\
\midrule
Qwen / LongVideoBench & .6223 & .6253 & ${+}0.30$ \\
Qwen / Video-MME & .6222 & .6178 & $-0.44$ \\
Qwen / LVBench & .4635 & .4674 & ${+}0.39$ \\
GPT-5-mini / LVBench & .4900 & .4926 & ${+}0.26$ ($p{=}.84$) \\
Uniform / LVB long pool & .5061 & .5092 & ${+}0.31$ ($p{=}.784$) \\
\midrule
\multicolumn{4}{l}{\emph{(b) Reinvestment at equal or lower token cost}} \\
Answerer / benchmark & $k{=}8$ full & $k{=}16$ compr. & $\Delta$ (pp) \\
\midrule
Qwen / LongVideoBench & .6223 & .6447 & ${+}2.24$ \\
Qwen / LVB long pool & .5820 & .6055 & ${+}2.36$ ($p{=}.0346$) \\
Qwen / Video-MME & .6222 & .6378 & ${+}1.56$ \\
Qwen / LVBench & .4635 & .4939 & ${+}3.04$ ($p{=}.0009$) \\
GPT-5-mini / LVBench & .4900 & .5139 & ${+}2.39$ ($p{=}.039$) \\
\bottomrule
\end{tabular}}
\caption{(a) Spatial compression at fixed timestamps. Qwen rows use the approximately half-budget arm; the GPT row uses a pixel-ratio control, and the uniform control an approximately 46\% realized per-frame budget. (b) Reinvestment; token parity is audited on the LongVideoBench bins only, and Video-MME and LVBench use the same resize rule without a separate audit.}
\label{tab:t2}
\end{table}

\section{Scorer-swap detail}
\label{app:scorer}

The 60\,s bin adds nothing either way (largest contrast ${+}3.49$, $p{=}.31$), as expected when eight frames already cover 22\% of a median 36-frame pool.
Discordant counts for the uniform contrasts are $b{=}69,c{=}37$ under LongCLIP and $b{=}65,c{=}29$ under SigLIP.

One reading must be resisted: OMP beats top-$k$ under SigLIP (${+}3.40$, $p{=}.049$) but not under LongCLIP (${+}2.43$, $p{=}.33$), which invites the conclusion that OMP suits SigLIP better.
% REV-2026-08-21 [R2] p corrected .73 -> .6752. The .73 came from a two-arm test
% (OMP vs uniform across scorers) in which uniform cancels; this sentence needs the
% four-arm difference-in-differences. Point estimates coincide only because top-k's
% marginal accuracy is identical under both scorers.
The interaction test rejects that: the difference of the two gains is 0.97 points ($b{=}48$, $c{=}43$, $p{=}.68$).
That estimate coincides with the OMP row of Table~\ref{tab:scorer} only because top-$k$ has identical marginal accuracy under both scorers; the two are different tests over different per-item patterns, which is why their $p$-values differ.
A significant result beside a non-significant one is not evidence that the two differ, and Appendix~\ref{app:qtype} records the same error made on question categories.


\section{Spatial allocation: controls and equivalence detail}
\label{app:compression}

On the pooled LongVideoBench long bins ($n{=}976$), the compressed OMP arm is 0.72 points \emph{higher} than full resolution, with a 90\% confidence interval of $[-0.60,+2.03]$.
That interval fits inside a $\pm3$-point margin (TOST $p{=}.0022$) but narrowly fails at $\pm2$ ($p{=}.054$), so $\pm3$ is the honest statement.
InternVL3-8B supports a stronger version under a more aggressive intervention: cutting 24 tiles to 8 moves pooled accuracy by $-0.10$ points with a 90\% interval of $[-1.66,+1.45]$, which clears a $\pm2$-point margin ($p{=}.022$).
The remaining LVBench and Video-MME differences in Table~\ref{tab:t2}a are descriptive stability checks; we do not present them as equivalence results.

% ACL-TRIM: figure `fig:claim-a` moved to appendix.


Compression being cheap does not mean the allocation \emph{shape} is irrelevant, and our controls separate the two.
The residual-proportional schedule does not beat a flat split at the same mean budget (Figure~\ref{fig:claim-a}), so OMP's ranking---useful as it is for choosing frames---did not improve accuracy detectably when used to distribute pixels among them. We ran no equivalence test on that contrast, so this is a failure to detect a difference rather than evidence of none.
A reconstruction of LDDR's stage-2 Group-DPP importance does better, improving on the flat split by 1.84 points on the pooled long bins ($p{=}.0198$, 95\% CI $[+0.37,+3.32]$).
That contrast comes from our own same-environment reconstruction rather than the authors' released implementation, and its 3600\,s component is not individually significant, so we treat it as indicative.
The conclusion is therefore three-part and narrower than ``resolution does not matter'': coarse compression is close to free, a naive priority-proportional rule adds nothing, and a better-designed importance signal can add something.


\section{What the later picks actually add}
\label{app:mechanism}

% ACL-TRIM: condensed ~354 -> ~230 words. All statistics and the audit counts
% retained; the exposition is tightened. Full version in arxiv.tex.

OMP's residual trace does not behave the way its signal-processing origin predicts. The account below is a hypothesis consistent with that discrepancy, not a demonstrated mechanism: nothing here intervenes on residual geometry.
On both LongVideoBench long bins the residual norm is 0.972 of its starting value after one pick and 0.967 after eight, falling no further by sixteen: almost none of the query is ever reconstructed.
What does change is the residual's correlation with the selected frame, dropping from .233 at pick one to .004 at pick eight and roughly zero by sixteen, so the query-directed signal is exhausted early.
Yet those later frames keep a cosine near .20 with the original query.

The text query sits almost orthogonal to the \emph{dominant directions} of the frame-embedding cone---the modality gap of contrastive image--text encoders~\citep{liang2022modality}, visible in Figure~\ref{fig:residual}b---so with an effectively low-rank frame dictionary, eight greedy picks reach only those directions and OMP cannot reduce the residual much whatever it selects.
What it can do is avoid repeating itself: late picks add little new query direction while adding chances for the answerer to encounter a relevant moment, which is what the reinvestment result rewards.

Coverage helps only while it stays relevant.
An exploratory inspection of 93 OMP failures at $k{=}8$ across the two long bins found one dominant pattern---off-topic novelty, in 20 of 41 inspected 600\,s cases and 35 of 52 inspected 3600\,s cases, where picks went to visually distinctive but question-irrelevant material.
The mirror-image failure also appeared: on cross-scene tracking questions, dense top-$k$ clustered on one relevant moment and missed another required scene.
Category-level $k{=}8$ versus $k{=}16$ results are mixed, so this supports a mechanism hypothesis rather than a rule about question types, but it points at a design target: a relevance floor, so spread is explored within plausible evidence regions.


\section{FOCUS$^\star$ schedule replay}
\label{app:focus}

Every row in Table~\ref{tab:t1}, including FOCUS$^\star$, consumes the same dense 1\,fps LongCLIP stem-score vector.
The original FOCUS method instead treats scoring itself as a budgeted online process, in which clip-arm pulls purchase ITM evidence~\citep{zhu2025focus}.
Our replay preserves the temporal clip schedule and default hyperparameters but reveals already-computed LongCLIP scores to the bandit.
It therefore tests the schedule under a matched scorer, and says nothing about the efficiency or accuracy of the full FOCUS pipeline.

The replay uses 16\,s clips, three coarse pulls per clip, $\alpha{=}0.25$, a pull budget equal to half the candidate pool, and approximately $k/4$ surviving clips.
Unobserved rewards are interpolated from the nearest observed frame, and the random seed is fixed per item.
This boundary is why the row carries a star, and why we draw no conclusion from it about FOCUS under its native ITM protocol.

\section{Question categories do not define a stable selector regime}
\label{app:qtype}

LongVideoBench provides question-category tags, which invite the hypothesis that temporally referred questions need different frame selection from the rest.
We test it by comparing OMP with uniform sampling on the temporally referred categories (the \texttt{T*} group) against all remaining categories, without relabeling the data.

\begin{table}[h]
\centering
\small
\resizebox{\columnwidth}{!}{%
\begin{tabular}{llrccc}
\toprule
Bin & Slice & $n$ & Uniform & OMP & $\Delta$ (McNemar) \\
\midrule
600\,s & temporal \texttt{T*} & 148 & .5608 & .5811 & ${+}2.03$, $p{=}.74$ \\
600\,s & other & 264 & .5492 & .6591 & ${+}10.98$, $p{=}5.2\times10^{-4}$ \\
3600\,s & temporal \texttt{T*} & 222 & .4369 & .5270 & ${+}9.01$, $p{=}.013$ \\
3600\,s & other & 342 & .4942 & .5585 & ${+}6.43$, $p{=}.021$ \\
\bottomrule
\end{tabular}}
\caption{OMP minus uniform on official LongVideoBench categories. Percentage-point gains differ by duration, and the category interaction is not significant in either long bin.}
\label{tab:qtype}
\end{table}

The temporal slice looks resistant to selection at 600\,s and responsive at 3600\,s, which is the shape of a result that would be easy to over-interpret.
It is not one.
The category interaction is not significant in either bin, and which slice ``wins'' reverses between them.
Many \texttt{T*} stems refer to subtitles or speech while our selector is vision-only and subtitles are disabled, which offers one plausible source of the instability.
Two significant-versus-non-significant subgroup tests are not an interaction, so we infer no question-type router from these results.

\section{Selector variants}
\label{app:negatives}

\begin{figure}[t]
\centering
\includegraphics[width=\columnwidth]{fig4_negative_sweep.pdf}
\caption{Coverage-gated selector variants against OMP at $k{=}8$ on the 600 and 3600\,s LongVideoBench bins. The shaded $\pm1.5$-point band is a visual aid, not a statistical criterion.}
\label{fig:negatives}
\end{figure}

This sweep changes only the subset rule, on the same LongCLIP cache.
MMR underperforms; query-blind DPP ($\beta{=}0$) collapses toward uniform sampling, which is what a diversity term carrying no query signal would look like in this one configuration; and query-weighted DPP, residual floors, and partial orthogonalization all sit in a band around OMP.
No tested DPP variant significantly exceeds OMP in either long bin.

This sweep ran in the secondary compute environment described in \S\ref{sec:limitations}; its arms are internally paired against each other and are never compared against Tables~\ref{tab:t1} and~\ref{tab:t2}.

These nulls do not show that all diversity objectives are equivalent.
They show something more specific and more useful: within this scorer geometry and budget range, moving around inside the relevance--diversity family produces less change than replacing uniform sampling or reallocating pixels into frames.
That is consistent with the residual analysis in \S\ref{sec:mechanism}---if the query-directed signal is exhausted after roughly five picks, then rules that differ only in how they trade relevance against diversity have little left to differ about.

\section{Exploratory failure audit}
\label{app:failure}

Two independent model-assisted visual passes inspected OMP failures at $k{=}8$ through a frame viewer: 41 cases spanning 16 categories in the 600\,s bin, and 52 cases spanning 17 categories in the 3600\,s bin.
Cases were chosen to cover categories and to include both OMP-only and shared failures; the sample is not random.
Reviewers compared the selected timestamps, the rendered frames, the question, the gold answer, and the model's prediction.
Primary tags were missed evidence, temporal offset, answerer error despite adequate frames, and questions not answerable from visual frames at all.

The most frequent tag in both passes was visually distinctive but irrelevant content consuming several of the eight picks---unrelated chapters in compilations, separate storylines in long narrative video, title and subscribe cards, and dark transitions. Because the sample is failure-conditioned, category-balanced by construction, model-assisted, and scored without blinding, inter-rater agreement, or a matched non-OMP control, we cannot call this pattern dominant, OMP-specific, or prevalent.
The opposite failure showed up on sequence-of-scenes and cross-scene tracking questions, where dense top-$k$ sampling covered one relevant moment but omitted another required scene.
The audit therefore motivates relevance-constrained diversity as a design direction.
It cannot estimate population failure rates, and it cannot establish that adding frames benefits any particular question category.

\end{document}
